\documentclass[10pt]{article}

\usepackage[utf8]{inputenc}
\usepackage[T1]{fontenc}
\usepackage[a4paper, left=2.54cm, right=2.54cm, top=3cm, bottom=3cm]{geometry}
\usepackage[spanish]{babel}
\usepackage{amsmath, amssymb}
\usepackage{multicol}
\usepackage{enumitem}
\usepackage{fancyhdr}
\usepackage{titlesec}

% ─────────────────────────────────────────────────────────────
% Configuración de página
% ─────────────────────────────────────────────────────────────

\pagestyle{fancy}
\fancyhf{}

\fancyhead[L]{Academia Prepolicial - CTA}
\fancyhead[R]{\today}

\fancyfoot[C]{\thepage}

\renewcommand{\headrulewidth}{0.8pt}

% ─────────────────────────────────────────────────────────────
% Configuración de títulos y espaciado
% ─────────────────────────────────────────────────────────────

\titleformat{\section}[block]
{\large\bfseries\centering}
{\thesection}{1em}{}

\setlength{\columnsep}{1cm}
\setlength{\parskip}{0.8em}
\setlength{\parindent}{0em}

% ─────────────────────────────────────────────────────────────

\begin{document}
	
	\begin{center}
		{\Large\bfseries Hoja de Ejercicios}\\[0.2cm]
		{\large Ciencia, Tecnología y Ambiente}\\[0.1cm]
		{\normalsize Materia, Energía y el Universo}
	\end{center}
	
	\section*{Ejercicios}
	
	\begin{multicols}{2}
		
		\begin{enumerate}[label=\textbf{\arabic*.}, itemsep=0.45cm]
			
			\item ¿Qué es la materia y cuáles son sus características principales?
			
			\item Mencione cuatro propiedades generales de la materia.
			
			\item ¿Cuál es la diferencia entre masa y peso?
			
			\item Calcule la densidad de un cuerpo de masa
			$18\,\mathrm{kg}$ y volumen $3\,\mathrm{m}^3$.
			
			\item ¿Qué propiedad de la materia permite que un metal pueda convertirse en hilos?
			
			\item ¿Qué es la elasticidad?
			
			\item Nombre los cuatro estados de la materia.
			
			\item ¿Cómo se llama el cambio de estado de líquido a gas?
			
			\item ¿Qué características posee el estado gaseoso?
			
			\item Un bloque tiene una densidad de
			$4\,000\,\mathrm{kg/m}^3$ y un volumen de
			$2\,\mathrm{m}^3$. Halle su masa.
			
			\item Defina energía.
			
			\item ¿Qué establece la Ley de Conservación de la Energía?
			
			\item Calcule la energía cinética de un cuerpo de
			$8\,\mathrm{kg}$ que se mueve a
			$5\,\mathrm{m/s}$.
			
			\[
			E_k = \frac{1}{2}mv^2
			\]
			
			\item Una caja de $10\,\mathrm{kg}$ se encuentra a una altura de $6\,\mathrm{m}$. Halle su energía potencial gravitacional considerando
			$g=9{,}8\,\mathrm{m/s^2}$.
			
			\[
			E_p = mgh
			\]
			
			\item Mencione tres tipos de energía.
			
			\item ¿Qué diferencia existe entre fuentes de energía renovables y no renovables?
			
			\item ¿Cuál es la principal fuente de generación eléctrica del Perú?
			
			\item ¿Qué es el Sistema Solar?
			
			\item Nombre los ocho planetas del Sistema Solar en orden desde el Sol.
			
			\item ¿Cuál es el planeta más grande del Sistema Solar?
			
			\item ¿Qué movimiento terrestre produce el día y la noche?
			
			\item ¿Qué movimiento terrestre origina las estaciones del año?
			
			\item Nombre las capas internas de la Tierra.
			
			\item ¿Qué es un mineral?
			
			\item ¿Cuál es la diferencia entre roca y mineral?
			
			\item Mencione un ejemplo de roca ígnea, sedimentaria y metamórfica.
			
			\item ¿Qué sostiene la teoría del Big Bang?
			
			\item ¿Cuál es la edad aproximada del Universo?
			
			\item ¿Qué evidencia demuestra la expansión del Universo?
			
			\item ¿Qué porcentaje del Universo corresponde a materia ordinaria?
			
		\end{enumerate}
		
	\end{multicols}
	
\end{document}